\documentclass[12pt]{article}
\usepackage[utf8]{inputenc}
\usepackage[T1]{fontenc}
\usepackage{hyperref}
\title{Aro Language Resources for K-12 Curriculum Development}
\author{Compiled by Claude Sonnet 5 (claude-sonnet-5) for Andr\'as Kornai}
\date{}
\begin{document}
\maketitle

\section{Introduction}
Aro (ISO 639-3 tei) is a Torricelli language of Papua New Guinea with about 500 speakers. It belongs to the Kombio branch and is spoken in the Torricelli Mountains, Sandaun Province. The one substantial resource found is a recent documentation project. Nothing for schools turned up.

\section{Primary Educational Resources}
No primers, textbooks or school programmes were found.

\section{Linguistic and Reference Materials}
\begin{itemize}
\item Aro language, Wikipedia overview: \url{https://en.wikipedia.org/wiki/Aro_language}
\item Documentation of Aro [tei], a Torricelli language of Papua New Guinea (Culture in Crisis project page): \url{https://cultureincrisis.org/projects/documentation-of-aro-tei-a-torricelli-language-of-papua-newguinea}
\end{itemize}

\section{Religious and Community Resources}
None specific to Aro were found. The Joshua Project pages that appeared concern ``Torricelli, Lou'', a different people group, and are not used here.

\section{Digital and Audio Resources}
\begin{itemize}
\item ELDP archive collection ``Documentation of Aro [tei]'' by Andrey Drinfeld. It holds video recordings of narratives, songs and conversations annotated in ELAN: \url{https://eldp.access.preservica.com/dk0613/}
\end{itemize}

\section{Curriculum Development Context}
The language is not known to be taught in schools. Adults are the only first-language users, so any curriculum would need to support intergenerational transmission. The ELDP recordings are the most useful starting material for community-led work.

\end{document}